\section{DATA as a distinguished case}
DATA 3.0 is a rich instance of the general framework. Its complete state includes object theory, target, metatheory, typed verified bank, multiple search frontiers, untrusted proposals, history, learned parameters, resource allocation, candidate/accepted certificates, settlement status, and a monotone round counter \cite{data3}. Its active subsystems are modeled as transformations whose compositions generate a transformation semigroup.

The general theory above interprets DATA in four simultaneous ways:
\begin{enumerate}
\item the instantaneous DATA process is a controlled transition system or IFS-like semigroup action;
\item its cumulative verified history is a SIC;
\item the $P$ route specializes to an ATP, while the combined $P/R/I$ routes form an ALD and the nested dependence/independence route forms an AMLD;
\item the resulting theorem/settlement graph is a rationally weighted maze once exact rational transaction costs are declared.
\end{enumerate}

The DATA-specific intelligence layer can therefore be understood as a collection of maze-navigation policies. Professor scores frontiers; Counselor challenges representations; Picard supplies feedback; Creativity and Dreamer generate candidate moves; Horizon certifies bounded regions; the verified bank adds reusable macro structure. None of these changes the verifier's semantics.

\section{A search funnel for computational accessibility}
The synthesis of the Horizon and Depths programs suggests the following funnel:
\[
\begin{array}{c}
\text{raw staged search dynamics}\\[2pt]
\downarrow\\[2pt]
\text{cumulative verified-history SIC}\\[2pt]
\downarrow\\[2pt]
\text{novelty / quotient / DAG reduction}\\[2pt]
\downarrow\\[2pt]
\text{conservative compilation and shared verified landmarks}\\[2pt]
\downarrow\\[2pt]
\text{verified incumbent bound }B\\[2pt]
\downarrow\\[2pt]
\text{rational branch-and-bound / Dijkstra / A* / exact graph methods}\\[2pt]
\downarrow\\[2pt]
\text{learned ranking, portfolios, reconnaissance, and resurvey}\\[2pt]
\downarrow\\[2pt]
\text{certificate-bearing fixed point.}
\end{array}
\]
The strongest claim is not that the fixed point becomes easy. It is that each justified transformation can remove a class of states or routes that do not need to be searched for the declared objective.

\section{Research conjectures and implementation questions}
\begin{conjecture}[Verified-history compression]
On broad formal-mathematics families, searching a canonicalized graph of cumulative verified histories requires asymptotically or practically fewer expansions than searching raw transient histories under matched correctness constraints.
\end{conjecture}

\begin{conjecture}[Compiled-landmark acceleration]
A verified shared bank plus conservative macro compilation reduces rational weighted hitting time on families with repeated substructure, after retrieval and verification overhead are charged exactly.
\end{conjecture}

\begin{conjecture}[Admissible learned lower bounds]
For some nontrivial theorem families, a learned function can be calibrated into a valid rational lower bound $h(v)$ often enough to improve branch-and-bound or A* search while preserving exact minimum-cost guarantees.
\end{conjecture}

\begin{conjecture}[Reflective search improvement]
A meta-ATP can prove bounded or conditional facts about a frozen source ATP/ALD/AMLD that, when deposited into the verified meta-bank, measurably reduce later settlement cost without altering the verifier or certificate semantics.
\end{conjecture}

Key implementation questions include whether cumulative-history canonicalization is cheap enough to pay for itself; how to define rational edge costs that remain meaningful across machines; how much of proof-state equivalence can be checked exactly; whether macro compilation should affect search cost, native proof cost, or both; and which meta-theorems are strong enough to improve scheduling rather than merely restating architecture.

\section{Conclusion}
The central proposal is that theorem proving and conjecture settlement can be studied as exact rationally weighted maze solving over verifier-backed state graphs. The internal search process may be nonmonotone, stochastic, learned, or self-reflective, while its cumulative history of verified information forms a monotone SIC. Under additional target-search conditions that SIC is an ATP; with typed proof/refutation/independence certificates it becomes an ALD; one logical level higher it becomes an AMLD.

The graph-theoretic perspective matters because it imports a large body of route-finding ideas subject to explicit hypotheses. For finite nonnegative rational weights, Dijkstra-style minimum-cost search is exact. With admissible lower bounds, A*-style and branch-and-bound pruning can reduce the region below an incumbent. Novelty DAGs eliminate verified-state cycles, conservative compilation replaces repeated verified paths by macro edges, and shared lemma banks turn repeated proof structure into reusable landmarks. These reductions do not solve undecidability or guarantee practical access to every fixed point. They narrow the ocean in which the machine must search.

The further reflective possibility is that the induced ATP can itself reason about a coded source ATP, ALD, AMLD, or controller. If such meta-theorems can be verified and fed back into search, the machine does not merely learn statistically from its trajectories; it can sometimes use proved facts about its own frozen search architecture. DATA is a distinguished laboratory for this general theory, not its definition.
